% FILE: 04_implementation_surface.tex
% Project: ggen — Governance Generation Engine

\section{Implementation Surface}
\label{sec:ggen-implementation}

\subsection{Source Files}
\label{sec:ggen-implementation-source}

The ggen implementation surface consists of specification files rather than compiled source code.
The \texttt{src/} directory exists but contains no Rust source files, confirming that ggen
operates as a template-and-query engine whose own implementation is external to this repository.
The authoritative source files are:

\textbf{Manifest:} \texttt{ggen.toml} (112 lines) defines the project configuration, four
active generation rules, the inference pass-through rule, query endpoint configuration
(\texttt{in-memory}, SPARQL 1.1), template engine configuration (Tera, \texttt{autoescape = true},
\texttt{strict\_variables = true}), and output configuration (\texttt{blake3},
\texttt{cryptographic-chain}).

\textbf{SPARQL Queries} (4 files in \texttt{queries/}):
\begin{itemize}
  \item \texttt{extract-board-claims.rq} --- 70 lines; SELECT with DISTINCT over 14 projected
    variables; enforces fitness $\geq 0.95$, precision $\geq 0.90$, and log format
    \texttt{ocel:2.0} or \texttt{xes:1849-2016} via FILTER.
  \item \texttt{extract-diligence-claims.rq} --- extends the board query with synergy category,
    debt remediation hours, risk severity, and replay trace deviations.
  \item \texttt{extract-lifecycle-governance.rq} --- extracts MAPE-K monitor, analyze, plan,
    execute, and knowledge rules per lifecycle state.
  \item \texttt{select-witness-projections.rq} --- extracts witness marker positions and lattice
    properties for Rust enum manufacturing.
\end{itemize}

\textbf{Governance Rules} (5 files in \texttt{rules/}): \texttt{feature-law.yaml} (542 lines)
defines six Cargo features with compliance gates; \texttt{ts-projection-law.yaml} defines
monomorphization rules, brand token rules, and eight verification checkpoints;
\texttt{wasm-boundary-law.yaml} defines ABI-safe type boundaries; \texttt{component-boundary-law.yaml}
defines WIT world segregation; \texttt{graduation-law.yaml} defines graduation readiness
conditions.

\textbf{Audit Scripts} (7 files in \texttt{audits/}, suffix \texttt{.sh.ggen}):
\texttt{audit-feature-law.sh.ggen} (6 gates), \texttt{audit-ts-projection-surface.sh.ggen},
\texttt{audit-ts-monomorphization.sh.ggen} (6 gates), \texttt{audit-ts-brand-tokens.sh.ggen}
(6 gates), \texttt{audit-ts-enum-tagging.sh.ggen} (8 gates),
\texttt{audit-no-engine-in-wasm-feature.sh.ggen}, and \texttt{audit-component-boundary.sh.ggen},
totaling 42 or more named gates.

\subsection{Commands}
\label{sec:ggen-implementation-commands}

The declared interface for invoking ggen is documented in \texttt{README.md} as a command-line
tool with the following invocation patterns:

\begin{itemize}
  \item \texttt{ggen --config ggen.toml --execute-all} --- execute all generation rules in the
    manifest sequentially, sealing each output with a BLAKE3 receipt.
  \item \texttt{ggen --config ggen.toml --rule <name>} --- execute a single named generation
    rule; e.g., \texttt{--rule blue-river-orchestrator} or \texttt{--rule wasm4pm-compat-witness-markers}.
  \item \texttt{ggen --config ggen.toml --verify-all} --- verify all receipts by re-running each
    originating query and comparing the resulting fitness score to the slide assertion.
\end{itemize}

The cargo-make integration targets documented in \texttt{GGEN\_MANUFACTURING\_SUMMARY.md} are:
\texttt{cargo make ggen-feature-plan}, \texttt{cargo make ggen-ts-projection},
\texttt{cargo make ggen-wasm-boundary}, \texttt{cargo make ggen-wit-world}, and
\texttt{cargo make audit-ggen-all}. \textit{(FUTURE WORK: the ggen binary does not exist as a
compiled artifact in this repository; whether these cargo-make targets invoke an external ggen
binary, a cargo plugin, or were executed manually is not conclusively documented in the available
artifacts.)}

\subsection{Data and Ontology Surfaces}
\label{sec:ggen-implementation-data}

\textbf{Ontology Extensions:} \texttt{ontology-extensions.ttl} defines the \texttt{ma:}
namespace (\texttt{https://process.intelligence/ma/}) with the \texttt{ma:BoardClaim} class
hierarchy and properties \texttt{ma:backedBy}, \texttt{ma:supportedBy}, \texttt{ma:evidencedBy},
\texttt{ma:quantifies}, and \texttt{ma:hasRemediationPath}. The \texttt{lifecycle:} namespace
defines \texttt{lifecycle:ProcessState} as a disjoint union of seven state classes with properties
for monitor, analyze, plan, execute, and knowledge rules. The \texttt{compat:} namespace defines
\texttt{compat:Evidence<T, State, W>} as a typed container for process evidence.

\textbf{wasm4pm-compat Ontology:} \texttt{wasm4pm-compat.ttl} contains the RDF extensions
specific to the wasm4pm-compat type law, including witness lattice positions and WASM ABI
type definitions.

\textbf{Query Endpoint:} The \texttt{[query]} section of \texttt{ggen.toml} specifies
\texttt{endpoint = "in-memory"} with \texttt{format = "sparql-11"}. The ontology namespace
prefixes are declared for \texttt{ma:}, \texttt{lifecycle:}, \texttt{wasm4pm:}, and
\texttt{compat:}. \textit{(INTERPRETATION: the in-memory endpoint implies that ggen loads the
ontology TTL files into a transient RDF store at manufacturing time rather than connecting to a
persistent SPARQL endpoint. Whether the conformance verdicts required by
\texttt{extract-board-claims.rq} were pre-populated in the TTL files or were absent from the
in-memory store during manufacturing is an open question documented in section
\ref{sec:ggen-open-questions}.)}

\textbf{Audit Log:} \texttt{audit.json} records \texttt{ggen\_version: 26.5.21},
\texttt{generated\_at: 2026-06-01T20:38:53.643973+00:00}, and \texttt{validation\_passed: true}.
The \texttt{inputs} object records empty hash maps for manifest, ontology, template, and query
hashes, and the \texttt{pipeline} and \texttt{outputs} arrays are empty, indicating the audit
record was committed before full pipeline execution rather than after it.

\subsection{Templates and Rendered Artifacts}
\label{sec:ggen-implementation-templates}

\textbf{\texttt{ma-deck.tera}:} Renders a seven-slide PowerPoint JSON structure: Title Slide,
Executive Summary, one Claim Detail Slide per board claim (with fitness, precision, receipt hash,
and event log format), Operational Debt Slide, Synergy Waterfall Slide, Conformance Model
Summary, and Board Sign-Off Slide with five admissibility declarations and a UUID timestamp.

\textbf{\texttt{ma-diligence.tera}:} Renders a seven-sheet Excel workbook JSON structure:
Executive\_Summary, Synergy\_Claims, Operational\_Debt, Integration\_Risks, Replay\_Traces,
Activity\_Impact, and Governance. Every row in every sheet is traceable to a SPARQL query result
and a conformance verdict.

\textbf{\texttt{blue-river.tera}:} Renders Rust source code for the \texttt{BlueRiverOrchestrator}
struct. The template iterates over the \texttt{states} variable (populated from
\texttt{extract-lifecycle-governance.rq} results) to manufacture the \texttt{LifecycleState}
enum, \texttt{StateDefinition} struct, \texttt{TransitionGuard}, \texttt{StateTransition},
\texttt{MonitorRule}, \texttt{AnalyzeRule}, \texttt{PlanRule}, \texttt{ExecuteAction}, and
\texttt{KnowledgeAsset} types. The first manufactured output of this template was sealed as
output hash \texttt{4e341f5a\ldots} in the chain-root receipt.

\textbf{\texttt{witness-marker.tera}:} Renders the \texttt{WitnessMarker} Rust enum with eight
named variants and an embedded test module with proofs of join-semilattice monotonicity (idempotency,
absorption, commutativity), partial order (reflexivity, transitivity), and two additional
algebraic invariants. The rendered artifact is
\texttt{../sources/wasm4pm-compat/manufactured/witnesses/witness-markers-20260601.rs}.
